Multi-agent systems built on large language models (LLMs) increasingly develop software, yet team size and the division of work are usually fixed by hand, proposed by a model, or found by costly search, and every extra agent adds tokens and coordination; we ask whether, given the design of the software, the work should be split at all, and into how many agents. We convert a Unified Modeling Language (UML) class diagram into a weighted graph of duties and select, before any code exists, the average-linkage team that maximizes modularization quality minus a per-agent cost, which may be a single agent; we also derive the expected dependency of equal-size random teams in closed form. On all \NProjects DevBench projects, the criterion recommends one agent for \NSingle and teams of \KMin to \KMax agents for \NMulti, with \CutRedMin to \CutRedMax percent less cross-agent dependency than random teams of equal sizes. In pre-registered experiments on RealBench with two model families, the criterion's teams used 27\% and 41\% fewer tokens than random teams of the same sizes, with no significant difference in tests passed, and with the first model avoided the collapse of one agent per class, which passed no test on 7 of 9 projects. Design-time sizing is thus a cheap, reproducible way to reduce the token cost of multi-agent software development and to avoid over-split teams, while a correctness advantage was not found. All code, prompts and cached responses are released.
